\chapter*{Prólogo}
\markboth{Prólogo}{Prólogo}
\addcontentsline{toc}{chapter}{Prólogo}

Las matemáticas nacieron como una herramienta para describir el mundo. Durante siglos, la geometría griega, el álgebra y posteriormente el cálculo permitieron resolver problemas de enorme complejidad con resultados sorprendentemente precisos. Sin embargo, ese éxito ocultaba una dificultad fundamental: muchos de los argumentos utilizados carecían de una justificación rigurosa.

El cálculo infinitesimal desarrollado por Newton y Leibniz revolucionó la ciencia, pero se apoyaba en conceptos como los \emph{infinitésimos}, los \emph{límites} y las \emph{series infinitas}, cuya naturaleza no estaba claramente definida. Los matemáticos obtenían resultados correctos, aunque con frecuencia los demostraban mediante argumentos intuitivos o manipulaciones que hoy consideraríamos insuficientes.

A lo largo del siglo~XIX comenzaron a aparecer contraejemplos que mostraban que la intuición no siempre era un guía confiable. Existían funciones continuas no derivables, series que cambiaban de valor al reordenar sus términos, sucesiones cuyo comportamiento desafiaba las ideas tradicionales y paradojas derivadas de un uso poco preciso del infinito. Estos fenómenos revelaron que las matemáticas necesitaban una base mucho más sólida.

La respuesta fue el nacimiento del \emph{análisis matemático moderno}. Matemáticos como Cauchy, Bolzano, Weierstrass, Dedekind, Cantor, Borel, Lebesgue, Fréchet, Hausdorff, Banach y muchos otros desarrollaron un nuevo lenguaje basado en definiciones precisas y demostraciones rigurosas. Conceptos como la \emph{completitud}, la \emph{continuidad}, la \emph{compacidad}, la \emph{convergencia} o la \emph{topología} dejaron de ser ideas intuitivas para convertirse en estructuras matemáticas perfectamente definidas.

A partir de ese momento, los grandes teoremas del análisis adquirieron un papel fundamental. Cada uno de ellos establece las condiciones exactas bajo las cuales un procedimiento matemático es válido: cuándo una sucesión converge, cuándo una función alcanza un máximo, cuándo una familia de funciones posee una subsucesión convergente o cuándo una función puede aproximarse arbitrariamente bien por polinomios. Los teoremas no son únicamente resultados aislados; son las garantías que permiten trabajar con el infinito de manera segura.

Este libro sigue ese recorrido histórico y conceptual. Comienza construyendo los fundamentos lógicos y axiomáticos sobre los que descansa el análisis, desarrolla la teoría de espacios métricos y las propiedades esenciales de las funciones, y culmina con algunos de los resultados estructurales más importantes del análisis moderno. Finalmente, muestra cómo estas ideas constituyen la base de disciplinas como la teoría de la medida, el análisis funcional, la teoría espectral, las ecuaciones diferenciales, la topología y la probabilidad moderna.

Más que aprender técnicas de cálculo, el objetivo de esta obra es comprender \emph{por qué las matemáticas funcionan}, cuáles son los principios que garantizan la validez de sus argumentos y cómo el rigor permitió transformar el cálculo intuitivo en una de las teorías más profundas y sólidas de toda la matemática.

\vspace{1em}
\begin{flushright}
\textit{M.H. Alberto}
\end{flushright}
